\section{Introduction}
Building an atomic model of a protein from a cryo-EM density map is still largely manual. MICA \citep{gyawali2025mica} automates it by combining the map with structures predicted by AlphaFold~3 \citep[AF3;][]{abramson2024af3}, and its authors report a clear advantage over ModelAngelo \citep{jamali2024modelangelo} and EModelX(+AF) \citep{chen2024emodelx}, including a mean TM-score of 0.93 on a set of recently released maps. This report asks a narrower question that a reader can check: \emph{how well does MICA build protein models on maps that were released in 2026, and how does it compare with the same two baselines and with newer methods, when everything runs on free compute and every method gets the same inputs?}

The test set is made of EMDB maps and PDB models that were both released between 1~January and 12~October 2026 (the freeze date). Twenty entries were chosen by a rule that was fixed before any model was run (Section~\ref{sec:datasets}). Five methods were run on all twenty entries with default settings and no per-map tuning (Section~\ref{sec:methods}); every output was scored in one session with one set of tools against the deposited model (Sections~\ref{sec:procedure} and~\ref{sec:metrics}). Failed runs stay in every table, and a sensitivity analysis scores them as zero. With about twenty entries the report puts effect sizes and confidence intervals before $p$-values.

\paragraph{Findings in brief.} MICA built a model for \NModelMica{} of the \NEntries{} entries (median TM-score \TmMedMica{}). The other methods were equal, slightly higher or lower depending on the method and the resolution; only the difference to EMProt is clear on both statistical tests (Section~\ref{sec:paired}). Two deviations from MICA's own pipeline, forced by free compute, and the small sample limit what can be concluded (Section~\ref{sec:limits}).
